% \begin{figure*}[!ht]
%     \centering
%     \begin{tikzpicture}[scale=0.9, every node/.style={scale=0.9}]
%         \definecolor{poscolor}{RGB}{0,158,115}
%         \definecolor{negcolor}{RGB}{213,94,0}
%         \definecolor{gencolor}{RGB}{86,180,233}
%         \definecolor{codebg}{RGB}{250,250,250}
%         \definecolor{textgray}{gray}{0.15}

%         % Positive Classification
%         \node[draw, fill=poscolor!20, text width=5cm, align=center, font=\small\bfseries, rounded corners=2pt]
%             at (-7, 6.2) {Positive Classification ($a$, $p$)};
%         \node[draw, fill=codebg, text width=5.8cm, align=left, font=\scriptsize, rounded corners=2pt]
%             at (-7, 4.5) {
%             \textbf{User:} Are these code snippets semantically equivalent?\\
%             Code 1: \texttt{[Anchor]}\\
%             Code 2: \texttt{[Positive]}\\
%             \textbf{Assistant:} Yes, they are equivalent...
%         };

%         % Negative Classification
%         \node[draw, fill=negcolor!20, text width=5.8cm, align=center, font=\small\bfseries, rounded corners=2pt]
%             at (0, 6.2) {Negative Classification ($a$, $n$)};
%         \node[draw, fill=codebg, text width=5.8cm, align=left, font=\scriptsize, rounded corners=2pt]
%             at (0, 4.5) {
%             \textbf{User:} Are these code snippets semantically equivalent?\\
%             Code 1: \texttt{[Anchor]}\\
%             Code 2: \texttt{[Negative]}\\
%             \textbf{Assistant:} No, they differ functionally...
%         };

%         % Forward Generation
%         \node[draw, fill=gencolor!20, text width=5.8cm, align=center, font=\small\bfseries, rounded corners=2pt]
%             at (7, 6.2) {Obfuscation Task ($a \rightarrow p$)};
%         \node[draw, fill=codebg, text width=5.8cm, align=left, font=\scriptsize, rounded corners=2pt]
%             at (7, 4.5) {
%             \textbf{User:} Obfuscate this code by \texttt{[Ofuscation Type]}...\\
%             \texttt{[Anchor]}\\
%             \textbf{Assistant:} Here's the obfuscated version:\\
%             \texttt{[Positive]}
%         };

%         % Learning Objectives
%         \draw[->, thick, poscolor] (-7, 3.4) -- (-7, 3.0);
%         \node[font=\normalsize, poscolor] at (-7, 2.8) {$\mathcal{L}_{pos}$: Minimize semantic distance};

%         \draw[->, thick, negcolor] (0, 3.4) -- (0, 3.0);
%         \node[font=\normalsize, negcolor] at (0, 2.8) {$\mathcal{L}_{neg}$: Maximize semantic distance};

%         \draw[->, thick, gencolor] (7, 3.4) -- (7, 3.0);
%         \node[font=\normalsize, gencolor] at (7, 2.8) {$\mathcal{L}_{gen}$: Learn transformation};

%         % Combined Loss
%         \node[draw, fill=white, rounded corners=3pt, text width=6cm, align=left, font=\large] 
%             at (0, 2) {$\mathcal{L}_{\text{CFT}} = \lambda_1 \mathcal{L}_{pos} + \lambda_2 \mathcal{L}_{neg} + \mathcal{L}_{gen}$};
%     \end{tikzpicture}
%     \caption{Contrastive Fine-Tuning prompting strategy.
%     %The anchor ($a$) is paired with a semantically equivalent transformation ($p$) and a dissimilar variant ($n$). The model learns to minimize semantic distance for $(a, p)$, maximize it for $(a, n)$, and generate $p$ from $a$.
%     }
%     \label{fig:contrastive_prompting}
% \end{figure*}

%%%%%%%%%%%%%%%%%%%%%%%%%%%%%%%%%%%%%%%%%%%%%%%%%%%

% \begin{figure*}[!ht]
%     \centering
%     \begin{tikzpicture}[scale=0.9, every node/.style={scale=0.9}]
%         \definecolor{poscolor}{RGB}{0,158,115}
%         \definecolor{negcolor}{RGB}{213,94,0}
%         \definecolor{gencolor}{RGB}{86,180,233}
%         \definecolor{codebg}{RGB}{250,250,250}
%         \definecolor{textgray}{gray}{0.15}

%         % Positive Classification
%         \node[draw, fill=poscolor!20, text width=5cm, align=center, font=\small\bfseries, rounded corners=2pt]
%             at (-7, 5.2) {Positive Classification ($a$, $p$)};
%         \node[draw, fill=codebg, text width=5.8cm, align=left, font=\scriptsize, rounded corners=2pt]
%             at (-7, 4.2) {
%             \textbf{User:} Are these code snippets semantically equivalent?\\
%             Code 1: \texttt{[Anchor]}\\
%             Code 2: \texttt{[Positive]}\\
%             \textbf{Assistant:} Yes, they are equivalent...
%         };

%         % Negative Classification
%         \node[draw, fill=negcolor!20, text width=5.8cm, align=center, font=\small\bfseries, rounded corners=2pt]
%             at (0, 5.2) {Negative Classification ($a$, $n$)};
%         \node[draw, fill=codebg, text width=5.8cm, align=left, font=\scriptsize, rounded corners=2pt]
%             at (0, 4.2) {
%             \textbf{User:} Are these code snippets semantically equivalent?\\
%             Code 1: \texttt{[Anchor]}\\
%             Code 2: \texttt{[Negative]}\\
%             \textbf{Assistant:} No, they differ functionally...
%         };

%         % Forward Generation
%         \node[draw, fill=gencolor!20, text width=5.8cm, align=center, font=\small\bfseries, rounded corners=2pt]
%             at (7, 5.2) {Obfuscation Task ($a \rightarrow p$)};
%         \node[draw, fill=codebg, text width=5.8cm, align=left, font=\scriptsize, rounded corners=2pt]
%             at (7, 4.2) {
%             \textbf{User:} Obfuscate this code by \texttt{[Ofuscation Type]}...\\
%             \texttt{[Anchor]}\\
%             \textbf{Assistant:} Here's the obfuscated version:\\
%             \texttt{[Positive]}
%         };

%         % Learning Objectives
%         \draw[->, thick, poscolor] (-7, 3.3) -- (-7, 3.0);
%         \node[font=\normalsize, poscolor] at (-7, 2.6) {$\mathcal{L}_{pos}$: Minimize semantic distance};

%         \draw[->, thick, negcolor] (0, 3.3) -- (0, 3.0);
%         \node[font=\normalsize, negcolor] at (0, 2.6) {$\mathcal{L}_{neg}$: Maximize semantic distance};

%         \draw[->, thick, gencolor] (7, 3.3) -- (7, 3.0);
%         \node[font=\normalsize, gencolor] at (7, 2.6) {$\mathcal{L}_{gen}$: Learn transformation};

%         % Combined Loss
%         \node[draw, fill=white, rounded corners=3pt, text width=6cm, align=left, font=\large] 
%             at (0, 1.5) {$\mathcal{L}_{\text{CFT}} = \lambda_1 \mathcal{L}_{pos} + \lambda_2 \mathcal{L}_{neg} + \mathcal{L}_{gen}$};
%     \end{tikzpicture}
%     \caption{Contrastive Fine-Tuning prompting strategy.
%     %The anchor ($a$) is paired with a semantically equivalent transformation ($p$) and a dissimilar variant ($n$). The model learns to minimize semantic distance for $(a, p)$, maximize it for $(a, n)$, and generate $p$ from $a$.
%     }
%     \label{fig:contrastive_prompting}
% \end{figure*}

%%%%%%%%%%%%%%%%%%%%%%%%%

